\documentclass[11pt,a4paper]{article}

% ---------- Core typography and page layout ----------
\usepackage[T1]{fontenc}
\usepackage[utf8]{inputenc}
\usepackage{lmodern}
\usepackage{microtype}
\usepackage[a4paper,margin=2.6cm]{geometry}
\usepackage{setspace}
\setstretch{1.05}
\emergencystretch=2em

% ---------- Mathematics ----------
\usepackage{amsmath,amssymb,mathtools}
\usepackage{bm}

% ---------- Tables and lists ----------
\usepackage{booktabs,longtable,array}
\usepackage{enumitem}
\providecommand{\tightlist}{%
  \setlength{\itemsep}{0pt}\setlength{\parskip}{0pt}}

% ---------- Quotations and links ----------
\usepackage{csquotes}
\usepackage{xurl}
\usepackage[hidelinks]{hyperref}
\hypersetup{
  pdftitle={Organizational Accessibility: From Self-Maintenance to Evolvability},
  pdfsubject={A Dynamical Framework for Cumulative Emergence}
}

% ---------- Bibliography ----------
% Overleaf: compile with Biber.
\usepackage[backend=biber,style=authoryear,natbib=true,doi=true,url=false,isbn=false,maxcitenames=2,maxbibnames=99]{biblatex}
\addbibresource{organizational_accessibility_v4.bib}
\DeclareFieldFormat{doi}{\mkbibacro{DOI}\addcolon\space\href{https://doi.org/#1}{\nolinkurl{#1}}}

% ---------- Document metadata ----------
\title{\textbf{Organizational Accessibility: From Self-Maintenance to Evolvability}\\[0.5em]
\large A Dynamical Framework for Cumulative Emergence}
\author{}
\date{Working paper --- 6 September 2026\\\small Scientific master v4}

\begin{document}
\maketitle

\begin{abstract}
Why can organization accumulate rather than repeatedly requiring
rediscovery from scratch? We propose that the relevant object is
\textbf{organizational accessibility}: the finite-horizon distribution
of organizational states that can follow from what already exists. The
claim is not merely that present state affects future state, which is
true of any dynamical system. It is that realized organization can
become part of the causal conditions determining which successor
organizations are effectively accessible.

We construct the account in four steps: a baseline self-maintenance
threshold followed by three distinct mechanisms. First, an open
production network can collectively replace the processes that
constitute it, yielding an organization reproduction number,
\(R_{\rm org}\). Below one, organization introduced at low abundance
decays; above one, it can establish. Second (\textbf{Mechanism I ---
hysteretic retention}), a seed-dependent feedback can separate
establishment from maintenance. In the worked model, a rare process
\(C\) cannot invade the resident \(A\)--\(B\) organization below
\(J_{\rm inv}=1.0667\), yet once established the expanded
\(A\)--\(B\)--\(C\) organization persists down to
\(J_{\rm fold}=0.5978\). At identical external conditions, realized
state and basin occupancy can therefore change \textbf{what can remain}.
Third (\textbf{Mechanism II --- inherited adaptive extension}),
reproduction with inheritance makes an already viable organization the
starting architecture for subsequent variation rather than requiring
viability to be reconstructed de novo. Continued reproduction also
supplies a finite exploratory opportunity before extinction. Near a
persistence boundary, a declining lineage can therefore combine many
remaining variant-generation trials with a comparatively small adaptive
displacement required to restore positive growth. This connects
organizational accessibility to evolutionary rescue and eco-evolutionary
feedback. Fourth (\textbf{Mechanism III --- evolvability}), the process
that generates candidate successors can itself vary: the distribution of
variants, their generation rate, or both may change and be
differentially retained.

We formalize the common object with a finite-horizon accessibility
kernel \(\mathcal A_\tau(S\mid\theta,z,E,Q)\), where \(\theta\) denotes
organizational architecture, \(z\) its current dynamical state, \(E\)
the external environment, and \(Q=(q,\nu)\) the variant-generating
process. For novel successors, exploratory opportunity

\[
\Omega_\tau
=
\mathbb E\!\left[\int_t^{t+\tau}N(s)\nu(s)\,ds\right]
\]

combines with the conditional variant distribution
\(q(\theta'\mid\theta)\) and establishment probability \(\pi_{\rm est}\)
to determine a successful-successor intensity. The three mechanisms are
mathematically distinct: Mechanism I changes effective persistence
accessibility through basin structure, Mechanism II changes the
inherited starting architecture and the amount of search available, and
Mechanism III changes \(Q\) itself.

The framework implies neither a force toward complexity nor
monotonically increasing accessibility. History may widen, narrow,
fragment, or collapse the set of viable successors. The general claim is
instead: \textbf{existence changes what existence can follow}.
Cumulative organization is the special case in which enough previous
viability is retained that it becomes infrastructure for subsequent
viability, so later viable organization need not repeatedly be
reconstructed from scratch.

\end{abstract}

\section{The problem: not survival of parts, but continuation of
organization}\label{the-problem-not-survival-of-parts-but-continuation-of-organization}

Organized systems persist despite continuous turnover. Molecules decay,
cells die, organisms replace cells, institutions replace members,
machines replace components, and computational processes terminate and
restart. Persistence therefore need not mean persistence of the same
material parts. It can instead mean persistence of a causal
organization: a network of processes whose activity repeatedly produces
the processes required for the network to continue.

This shifts the explanatory target. The first question is not merely why
a component survives, but how a collection of locally acting processes
can collectively cause its own continuation. Once such continuation is
possible, a second question appears:

\begin{quote}
\textbf{How can realized organization become part of the causal
conditions determining the accessibility of its own successors?}
\end{quote}

This formulation is deliberately stronger than the trivial statement
that present state affects future state, but weaker than claiming that
every realized state rewrites the underlying laws of transition. The
claim developed here is about \textbf{effective organizational
accessibility}. History can matter because an established state occupies
a different basin, because descendants inherit an already viable
architecture, because reproductive history changes the number of
exploratory trials available, because realized organization changes its
environment, or because the variant-generating process itself changes.

The argument proceeds in four steps. Sections 2--3 derive a baseline
reproduction threshold for self-maintaining organization. Sections 5--6
establish \textbf{Mechanism I --- hysteretic retention}: a newly
established feedback can alter the conditions required for its own
continued existence. Section 7 develops \textbf{Mechanism II ---
inherited adaptive extension}: once organization is reproduced,
successor search begins from inherited viable architecture and is
supplied by a finite, history-dependent exploratory opportunity. Section
8 develops \textbf{Mechanism III --- evolvability}: the process
generating candidate successors can itself vary and be retained.

The construction sits near several established traditions.
Autocatalytic-set theory studies networks in which reactions are
collectively sustained by catalysts produced within the network
\citep{Hordijk2010, Hordijk2012}. Jain and Krishna showed how
autocatalytic structure can emerge and reorganize in evolving catalytic
networks \citep{Jain1998, Jain2001}. Peng, Linderoth and Baum formalized
seed-dependent autocatalytic systems that can remain active after a
transient seed and can form hierarchical structures \citep{Peng2022}.
Evolutionary-rescue theory studies how populations facing environmental
deterioration can avoid extinction through adaptation before decline is
complete \citep{Bell2017, Ferriere2013}, while pathogen-emergence theory
has shown that subcritical transmission chains become longer as a
reproduction number approaches one, increasing cumulative opportunities
for adaptive change \citep{Antia2003, Woolhouse2005}.

The present contribution is narrower than a general theory of
complexity. We use an explicit resource balance to derive a reproduction
threshold, show how one added feedback produces a backward-bifurcation
geometry, distinguish that mechanism from inheritance-mediated adaptive
extension near a persistence boundary, and connect both to variation in
the process that generates successor organizations. The aim is not to
prove that complexity increases. It is to identify measurable ways in
which organizational history can condition what organization is
accessible next.

\subsection{Organizational
accessibility}\label{organizational-accessibility}

To keep architecture distinct from instantaneous state, let

\[
\theta=\text{organizational architecture},
\qquad
z=\text{instantaneous dynamical state}.
\]

For example, in the models below \(\theta\) may specify the available
process types and interaction structure, while \(z\) contains their
current abundances and the resource state. Let \(E\) denote the external
environment. Let

\[
Q=(q,\nu)
\]

denote the process generating organizational variants, where
\(q(\theta'\mid\theta)\) is the conditional distribution of candidate
successor architectures given a variation event and \(\nu\) is the rate
of such events per active realization.

For a measurable region \(S\) of the joint architecture--state space and
a finite horizon \(\tau\), define the
\textbf{organizational-accessibility kernel}

\[
\boxed{
\mathcal A_\tau(S\mid\theta,z,E,Q)
=
\Pr\!\left(
(\Theta_{t+\tau},Z_{t+\tau})\in S
\mid
\Theta_t=\theta,
Z_t=z,
E,Q
\right).
}
\tag{1}
\]

The exact state space and horizon must be specified in any real
application; \(E\) and \(Q\) may denote fixed values or specified
exogenous trajectories over that horizon. In a deterministic model, this
kernel may collapse onto one or several outcomes determined by
architecture, dynamical state, and basin occupancy. In a stochastic or
evolutionary system, it is a non-trivial distribution over successor
organizational states.

Equation (1) is a transition kernel at a specified horizon. If the
empirical question is instead whether the process \textbf{ever reaches}
a viable successor before time \(\tau\), the appropriate object is a
finite-horizon hitting probability. That distinction is useful because
persistence, first passage, and novel-successor generation need not have
the same mathematical form.

For novel-successor search, define \textbf{exploratory opportunity} as

\[
\boxed{
\Omega_\tau
=
\mathbb E\!\left[
\int_t^{t+\tau}N(s)\nu(s)\,ds
\right],
}
\tag{2}
\]

where \(N(s)\) is the number of active realizations. In a deterministic
or mean-field treatment, the expectation may be omitted. \(\Omega_\tau\)
is the expected number of relevant candidate-generation events over the
horizon; it is not the standing population and not an independent
primitive of the system, because it is generated by \((\theta,z,E,Q)\)
and by prior history.

A generated candidate still need not establish. Let

\[
\pi_{\rm est}(\theta'\mid\theta,z,E)\in[0,1]
\]

denote the probability that candidate architecture \(\theta'\)
establishes or otherwise enters the realized trajectory when generated
from current architecture \(\theta\) and state \(z\) in environment
\(E\). For a measurable architecture set \(U\subseteq\Theta\), a minimal
intensity of successfully established novel successors in \(U\) is then

\[
\boxed{
\Lambda_\tau(U\mid\theta,z,E,Q)
=
\Omega_\tau
\int_U
q(\theta'\mid\theta)
\pi_{\rm est}(\theta'\mid\theta,z,E)
\,d\theta'.
}
\tag{3}
\]

Under an independent Poisson approximation,

\[
\Pr(\text{at least one successful novel successor in }U)
\approx
1-e^{-\Lambda_\tau(U)}.
\tag{4}
\]

Equation (3) gives a useful conceptual decomposition:

\[
\boxed{
\text{successful-successor accessibility}
\sim
\text{opportunities}
\times
\text{what is generated}
\times
\text{what can establish}.
}
\]

The factorization is not claimed to be exact for every system. Its value
is that it separates quantities that can often be perturbed or estimated
independently.

The three mechanisms developed below enter this framework through
different causal channels:

\begin{itemize}
\tightlist
\item
  \textbf{Mechanism I --- hysteretic retention.} At fixed external
  conditions, realized dynamical state and basin occupancy can determine
  whether an already established organization remains. Realization
  changes \textbf{what can remain}.
\item
  \textbf{Mechanism II --- inherited adaptive extension.} Inheritance
  makes \(\theta\) the base point from which \(q(\theta'\mid\theta)\) is
  sampled, while reproductive history changes the \(N(t)\) contribution
  to \(\Omega_\tau\). Realization changes \textbf{where search begins
  and how much search occurs}.
\item
  \textbf{Mechanism III --- evolvability.} Organization can alter
  \(Q=(q,\nu)\) itself. Realization changes \textbf{how candidate
  successors are generated}.
\end{itemize}

The environment \(E\) remains conceptually distinct. In eco-evolutionary
systems, realized organization may also alter a future environment
\(E_{t+1}\), thereby changing later accessibility through an additional
feedback. The SARS-CoV-2 illustration in Section 7.3 uses this route,
but environmental modification is not required for the framework.

None of these mechanisms need increase complexity, adaptability, or
persistence. Accessibility may widen, narrow, shift, fragment, or
collapse. The most general statement is descriptive:

\[
\boxed{\text{Existence changes what existence can follow.}}
\]

Cumulative emergence is a more specific regime. It occurs when
historical changes in accessibility carry forward enough causal
structure for viability that later viable organization can be reached
without reconstructing prior viability de novo. In that case,

\[
\boxed{\text{previous viability can become infrastructure for subsequent viability}.}
\]

This motivates the working definition used throughout the paper:

\begin{quote}
A \textbf{cumulative organizational ratchet} occurs when realized
organization changes subsequent effective accessibility in a way that
carries forward enough previously achieved viability that later viable
organization need not be rediscovered de novo.
\end{quote}

\section{A resource-limited self-maintaining production
network}\label{a-resource-limited-self-maintaining-production-network}

Let \(x_i \geq 0\) denote the abundance of executing process type \(i\),
and let \(r \geq 0\) denote activated substrate. The nonnegative matrix
\(B\) contains local production relations: \(B_{ij}\) is the rate at
which process \(j\) produces process \(i\) per unit substrate and
producer abundance. Each new unit consumes one unit of substrate.
Existing processes are lost at rate \(d > 0\). The environment supplies
activated substrate at rate \(J\), while unused substrate is lost at
rate \(\ell r\). The model is

\[
\dot{x} = rBx - dx, \tag{5}
\] \[
\dot{r} = J - \ell r - r\mathbf{1}^T Bx. \tag{6}
\]

The same production term enters positively in the process balance and
negatively in the substrate balance. Hence this is a replacement model
rather than a model in which producers are immortal. Defining

\[
T = r + \sum_i x_i,
\]

we obtain

\[
\dot{T} = J - \ell r - d\sum_i x_i \leq J - \min(\ell, d)T.
\]

Thus total activated material remains bounded, and when \(J = 0\) the
active organization ultimately disappears. The organization is causally
closed only in the limited sense that its processes can regenerate one
another; energetically and materially it remains open.

This distinction is important. The model is not a derivation from full
chemical thermodynamics. A real molecular implementation would require
stoichiometric, kinetic and free-energy constraints to be specified
separately. Thermodynamic work on molecular replicators makes clear that
replication and selection carry energetic requirements rather than
escaping them \citep{Kolchinsky2025}.

\section{A reproduction number for
organization}\label{a-reproduction-number-for-organization}

With no active organization, the substrate converges to

\[
r_0 = \frac{J}{\ell}.
\]

Near this organization-free equilibrium, the process dynamics linearize
to

\[
\dot{x} = (r_0 B - dI)x.
\]

Let \(\rho(B)\) denote the spectral radius of \(B\). The leading growth
rate is

\[
s = r_0 \rho(B) - d,
\]

which motivates the dimensionless organization reproduction number

\[
R_{\rm org} = \frac{J\rho(B)}{\ell d}. \tag{7}
\]

If \(R_{\rm org} < 1\), sufficiently small organization decays. If
\(R_{\rm org} > 1\), it can grow from low abundance. For irreducible
\(B\), let \(p\) be the positive Perron eigenvector normalized by
\(\mathbf{1}^T p = 1\) and let \(\lambda = \rho(B)\). The positive
equilibrium is

\[
r^* = \frac{d}{\lambda}, \qquad x^* = Np, \qquad N = \frac{J - \ell d/\lambda}{d},
\]

which exists precisely when \(R_{\rm org} > 1\).

The same calculation can be written in next-generation-matrix form
\citep{VandenDriessche2002}. At the organization-free equilibrium,

\[
F = r_0 B, \qquad V = dI,
\]

so

\[
K = FV^{-1} = \frac{r_0}{d}B, \qquad \rho(K) = \frac{J\rho(B)}{\ell d} = R_{\rm org}.
\]

The vocabulary is borrowed from infectious-disease dynamics, but the
interpretation is generic: \(R_{\rm org}\) asks whether one generation
of executing processes collectively produces more than one replacement
generation before being lost.

For the two-process cycle

\[
B = \begin{pmatrix} 0 & \alpha \\ \beta & 0 \end{pmatrix},
\]

we have \(\rho(B) = \sqrt{\alpha\beta}\). Neither process need reproduce
itself directly. A closed return path across multiple process types can
be enough. If the return path is broken so that the matrix is acyclic,
the spectral radius falls to zero. This isolates the causal object of
interest: not activity alone, but a route from present activity back to
renewed productive capacity.

The positive equilibrium of the base model appears continuously at
\(R_{\rm org} = 1\):

\[
N = \frac{J}{d}\left(1 - \frac{1}{R_{\rm org}}\right).
\]

This is the expected forward transcritical geometry. In the base model,
establishment from rarity and loss of the positive equilibrium meet at
the same threshold. Section 6 shows how a seed-dependent feedback can
separate establishment from maintenance and thereby create hysteresis.
Section 7 asks a different question while retaining the same basic
threshold: how much exploratory opportunity remains, and how large an
adaptive change is required, when an established organization is driven
back toward and below \(R_{\rm org}=1\).

\section{An innovation must pay its opportunity
cost}\label{an-innovation-must-pay-its-opportunity-cost}

A new process is not automatically an improvement. To make this
explicit, begin with \(A \to B\) and \(B \to A\) at unit coefficient.
Let \(A\) divert a fraction \(u\) of its production capacity from \(B\)
to a new process \(C\), while \(C\) contributes back to production of
\(A\) with coefficient \(v\):

\[
B(u,v) = \begin{pmatrix} 0 & 1 & v \\ 1-u & 0 & 0 \\ u & 0 & 0 \end{pmatrix}, \qquad 0 < u < 1,\ v \geq 0.
\]

The dominant eigenvalue is

\[
\lambda(u,v) = \sqrt{1-u+uv} = \sqrt{1+u(v-1)}.
\]

Hence the new process lowers the external supply required for
maintenance only when \(v > 1\). If \(v < 1\), the diversion to \(C\)
reduces collective replacement capacity. The model therefore contains no
rule that rewards novelty, diversity, or complexity. An addition must
causally return enough productive benefit to offset its cost.

This distinction matters for the later argument. None of the mechanisms
below is generated by a preference for larger networks. Mechanism I
changes the persistence component of accessibility because a new
feedback loop alters the persistence geometry (Sections 5--6). Mechanism
II instead changes the inherited starting point and the number of
exploratory trials available before loss (Section 7). Mechanism III
changes the distribution from which candidate successors are generated
(Section 8).

\section{A seed-dependent extension: history enters the
state}\label{a-seed-dependent-extension-history-enters-the-state}

A fixed linear production matrix does not capture an innovation that
becomes self-maintaining only after it is present. We therefore consider
a nonlinear extension with reactions represented by

\[
R + B \to A + B, \qquad R + A \to A + B, \qquad R + A + C \to A + 2C, \qquad R + C \to A + C.
\]

The corresponding normalized equations are

\[
\dot a = rb + vrc - a, \tag{8}
\] \[
\dot b = ra - b, \tag{9}
\] \[
\dot c = urac - c, \tag{10}
\] \[
\dot r = J - r - (rb + vrc + ra + urac). \tag{11}
\]

All production consumes substrate, and \(\dot T = J - T\) for
\(T = r+a+b+c\).

The \(A\)--\(B\) subsystem can exist without \(C\). Process \(C\),
however, reproduces only when both \(A\) and \(C\) are already present,
while \(C\) contributes back to \(A\). A one-off seed can therefore
activate a loop that was dynamically inaccessible from \(c=0\).

For \(u=30\) and \(v=2\), numerical experiments give the following
controlled comparison. Start the \(A\)--\(B\) system at its equilibrium
with \(J=1.2\). At time \(t=50\), introduce \(10^{-4}\) of \(C\) while
subtracting the same amount from resource. At \(t=150\), lower supply to
\(J=0.9\). With the \(C \to A\) return path intact, \(A\), \(B\), and
\(C\) approach a positive equilibrium. If \(C\) is never seeded, the
organization disappears after the supply reduction. If \(C\) is seeded
but its return effect on \(A\) is removed (\(v=0\)), the organization
also disappears. Thus the persistent difference is carried by the newly
established feedback architecture, not by continuing external supply of
\(C\).

The seed need not be deliberately introduced. One may add a rare local
reaction \(R+A \to A+C\) and treat its occurrence stochastically. Once a
rare event creates \(C\), the ordinary deterministic dynamics can
amplify it. This does not explain why such a reaction must exist in
nature; it shows only that an external evaluator that recognizes and
installs the ``better'' organization is not logically required.

\section{Mechanism I --- hysteretic retention: changing what can
remain}\label{mechanism-i-hysteretic-retention-changing-what-can-remain}

The key result is that the threshold for establishing \(C\) is not the
same as the threshold for maintaining the \(C\)-containing state.

For \(c>0\), the equilibrium conditions of Eq. (11) imply

\[
a = \frac{1}{ur}, \qquad b = \frac{1}{u}, \qquad c = \frac{1-r^2}{uvr^2}, \qquad 0<r<1.
\]

The associated resource supply is

\[
J(r) = r + \frac{1/r + 1 + (1/r^2-1)/v}{u}. \tag{12}
\]

The fold of this positive equilibrium branch satisfies

\[
\frac{dJ}{dr}=0 \iff ur^3 - r - \frac{2}{v} = 0.
\]

For \(u=30\) and \(v=2\),

\[
r_{\rm fold} \approx 0.356236, \qquad J_{\rm fold} \approx 0.597806.
\]

Now examine invasion from the resident \(A\)--\(B\) state. For \(J>1\),
that equilibrium has

\[
r^*=1, \qquad a^*=b^*=\frac{J-1}{2}.
\]

At very small \(c\),

\[
\dot c \approx (ur^*a^*-1)c,
\]

so the type-reproduction number of \(C\) is

\[
R_C = ur^*a^* = u\frac{J-1}{2}.
\]

The invasion threshold \(R_C=1\) is therefore

\[
J_{\rm inv} = 1+\frac{2}{u} = 1.066667. \tag{13}
\]

This is the same logic used for type-specific invasion numbers in
heterogeneous epidemic systems \citep{Heesterbeek2007}.

The striking fact is

\[
J_{\rm fold} < 1 < J_{\rm inv}.
\]

The positive \(C\)-containing equilibrium branch extends backward from
the invasion point to much lower values of \(J\), turning at the fold.
This is the characteristic geometry called a backward bifurcation in
epidemiological models with sub-threshold endemic equilibria
\citep{KribsZaleta2000, VandenDriessche2002}.

The stable states can be summarized as follows:

\begin{longtable}[]{@{}lll@{}}
\toprule\noalign{}
Supply range & Stable unexpanded state & Stable expanded state \\
\midrule\noalign{}
\endhead
\bottomrule\noalign{}
\endlastfoot
\(J < J_{\rm fold}\) & organization-free & no \\
\(J_{\rm fold} < J < 1\) & organization-free & \(A\)--\(B\)--\(C\) \\
\(1 < J < J_{\rm inv}\) & \(A\)--\(B\) & \(A\)--\(B\)--\(C\) \\
\(J > J_{\rm inv}\) & \(A\)--\(B\) unstable to rare \(C\) &
\(A\)--\(B\)--\(C\) \\
\end{longtable}

Thus, at \(J=0.9\), an unexpanded system decays to the organization-free
state while an already established \(A\)--\(B\)--\(C\) organization
persists. At \(J=1.02\), a small \(C\) introduced into the \(A\)--\(B\)
state cannot invade, but an already established \(A\)--\(B\)--\(C\)
state can persist. In both cases the same current environment permits
different long-term outcomes depending on history. This is hysteresis.

In the accessibility formalism of Section 1.1, the effect belongs in the
\textbf{persistence component of \(\mathcal A\)}, not in the external
environment \(E\) itself. The environment can be numerically identical
in the two histories, yet the probability mass that \(\mathcal A\)
assigns to continued occupancy of the expanded organization differs
because the realized feedback architecture has altered the system's
persistence geometry.

\begin{quote}
An organizational innovation exhibits \textbf{Mechanism I --- hysteretic
retention} when the conditions required to establish it are stricter
than the conditions required to maintain it once established.
\end{quote}

The backward bifurcation is not synonymous with organizational
accessibility; it is one explicit dynamical mechanism by which
accessibility becomes history-dependent. The new loop changes its own
persistence conditions. Once it exists, simply reversing the
environmental change that permitted its establishment does not
necessarily undo it. This is stronger than saying that positive feedback
stabilizes a system. It says that realized organization changes the
retention component of future accessibility: \textbf{what can remain
depends on what has already been realized.}

\section{Mechanism II --- inherited adaptive extension: changing where
search
begins}\label{mechanism-ii-inherited-adaptive-extension-changing-where-search-begins}

Mechanism I changes what can remain by altering the retention geometry
of an already realized organization. Mechanism II is different. It does
not require bistability, a fold, or a newly closed feedback loop. It
arises whenever a reproducing organization generates descendants from
its current structure rather than from the organization-free state.

Two consequences follow.

First, inheritance changes the \textbf{base point} of subsequent search.
Variants are generated around an already viable architecture \(\theta\),
not around the organization-free architecture.

Second, reproduction changes the \textbf{number of draws} made from the
local variation process before the lineage disappears. This is the
exploratory-opportunity term \(\Omega\) introduced in Eq. (2).

Together these effects create \textbf{inheritance-mediated adaptive
extension}: what has already survived changes both the base point of
future variation and the amount of exploration available before
persistence is lost. During environmental deterioration, the same
ingredients can additionally produce evolutionary rescue.

\subsection{Exploratory opportunity near a reproduction
threshold}\label{exploratory-opportunity-near-a-reproduction-threshold}

Consider a simple branching process with mean reproduction number
\(R<1\). Let \(N_g\) be the expected number of active realizations in
generation \(g\), beginning from \(N_0\) founders. Then

\[
E[N_g]=N_0R^g.
\]

The expected \textbf{cumulative number of realizations} before
extinction is therefore

\[
E[C]
=
\sum_{g=0}^{\infty}E[N_g]
=
\frac{N_0}{1-R},
\qquad R<1.
\tag{14}
\]

For one founder, this reduces to the familiar \(1/(1-R)\) result.
Importantly, Eq. (14) is a statement about cumulative lineage size, not
the instantaneous standing population.

Now let \(\nu\) denote the expected number of relevant
candidate-variation events per realization. Under the simplest
independent approximation, the expected exploratory opportunity
generated before extinction is

\[
\boxed{
\Omega
\approx
\nu E[C]
=
\frac{\nu N_0}{1-R}.
}
\tag{15}
\]

Equation (15) is the corresponding deterministic/mean-field
approximation to Eq. (2). Thus, for fixed \(N_0\) and \(\nu\), a lineage
driven only slightly below its persistence threshold can still generate
many more candidate variants before disappearing than a lineage driven
far below it. For example, one founder gives \(E[C]=2\) at \(R=0.5\),
\(10\) at \(R=0.9\), and \(100\) at \(R=0.99\).

The divergence as \(R\to1^-\) is an idealized branching-process result.
Finite populations, finite observation windows, competition, resource
limitation, changing environments, and correlations among descendants
will truncate it. The relevant general quantity is not \(1/(1-R)\)
itself but \(\Omega\): the integrated number of opportunities available
for generating candidate successors before the current organization is
lost.

This also clarifies how pre-establishment and post-establishment
situations differ. Before establishment, a rare organization typically
has small \(N_0\), so even a favourable near-threshold value of \(R\)
may yield only a limited search. After a period of supracritical
reproduction, deterioration may begin from a much larger \(N_0\). The
combination of an inherited viable structure and a large initial
population can make \(\Omega\) orders of magnitude larger than during
the original emergence.

Standing variation present when deterioration begins provides an
additional source of candidates, but it should not be folded into
\(\Omega_\tau\). Exploratory opportunity counts \textbf{new
candidate-generation events}. Pre-existing candidate diversity
contributes separately to the successful-successor intensity. A useful
decomposition is

\[
\boxed{
\Lambda_\tau(U)
=
\Lambda_{\rm standing}(U)
+
\Omega_\tau
\int_U
q(\theta'\mid\theta)
\pi_{\rm est}(\theta'\mid\theta,z,E)
\,d\theta',
}
\tag{16}
\]

where \(\Lambda_{\rm standing}(U)\) denotes the expected
successful-establishment contribution, over the same horizon, of
candidate variants already present at the onset of deterioration.
Candidates generated after that onset belong only to the
\(\Omega_\tau q\) term, avoiding double counting. This keeps
\textbf{standing variation} conceptually distinct from \textbf{new
exploratory opportunity generated during decline}.

\subsection{Near threshold, the adaptive step required for rescue can be
small}\label{near-threshold-the-adaptive-step-required-for-rescue-can-be-small}

Now suppose an established organization experiences environmental
deterioration so that its resident reproduction number falls just below
one:

\[
R_{\rm res}=1-\varepsilon,
\qquad
0<\varepsilon\ll1.
\]

A variant with relative reproductive advantage \(s\) has approximately

\[
R_{\rm var}
=
(1-\varepsilon)(1+s).
\]

The variant grows when

\[
(1-\varepsilon)(1+s)>1,
\]

or

\[
s>
\frac{\varepsilon}{1-\varepsilon}
\approx
\varepsilon
\qquad
(\varepsilon\text{ small}).
\tag{17}
\]

A more general local form makes the geometry explicit when \(\theta\)
admits a differentiable local parameterization. Let \(R(\theta,E)\)
denote the continuation number of architecture \(\theta\) in environment
\(E\), with \(R(\theta,E)=1-\varepsilon\). For a nearby variant
\(\theta'=\theta+\delta\theta\),

\[
R(\theta+\delta\theta,E)
\approx
1-\varepsilon
+
\nabla_\theta R(\theta,E)\cdot\delta\theta.
\]

Positive growth is locally restored whenever

\[
\boxed{
\nabla_\theta R(\theta,E)\cdot\delta\theta>\varepsilon.
}
\tag{18}
\]

This identifies a geometric opportunity, not a guarantee. Accessible
variants may not include a favourable direction; the local gradient may
be small; adaptation may require crossing a fitness valley; or decline
may be too rapid for a suitable candidate to arise and establish.

The vicinity of the threshold is therefore special for two independent
reasons.

First, the adaptive displacement required to restore positive growth is
small.

Second, for fixed \(N_0\) and variant-generation rate \(\nu\),
exploratory opportunity \(\Omega_\tau\) is large because decline is
slow.

These two effects multiply in the novel-successor intensity. A
rescue-capable architecture region \(U_{\rm rescue}\subseteq\Theta\) is
more likely to be reached when

\[
\Lambda(U_{\rm rescue})
=
\Omega
\int_{U_{\rm rescue}}
q(\theta'\mid\theta)\,
\pi_{\rm est}(\theta'\mid\theta,z,E)\,d\theta'
\]

is large. Near threshold, \(\Omega\) increases while the set of variants
capable of restoring \(R>1\) can also expand because the required
advantage is smaller.

This is the logic of \textbf{evolutionary rescue}: decline does not
imply instantaneous extinction. It leaves a finite interval during which
already existing diversity and newly generated variants may contain or
produce a successor whose growth rate is positive in the changed
environment \citep{Bell2017}. Rescue is not guaranteed. It depends on
the initial abundance, the variation kernel \(q\), standing variation,
the rate of new variant production, and the establishment filter
\(\pi_{\rm est}\), as well as on how quickly and how far the environment
has deteriorated.

Schematically,

\[
\text{deterioration}
\rightarrow
R_{\rm res}<1
\rightarrow
\text{finite }\Omega_\tau
\rightarrow
\text{candidate successor}
\rightarrow
R_{\rm var}>1
\rightarrow
\text{renewed expansion}.
\]

Unlike Mechanism I, no bistable branch is required. Mechanism II is
therefore not a second kind of hysteresis. It is
\textbf{inheritance-mediated adaptive extension under finite exploratory
opportunity}.

\subsection{Illustration: SARS-CoV-2 and an eco-evolutionary
feedback}\label{illustration-sars-cov-2-and-an-eco-evolutionary-feedback}

Pathogen evolution makes this mechanism unusually intuitive.

At the start of the SARS-CoV-2 pandemic, the human population was
largely immunologically naive. Viral circulation then changed that
environment by generating infection-derived immunity, while vaccination
added further immune pressure. The ecological consequence of prior viral
success therefore altered the subsequent fitness landscape.

As population immunity accumulated, lineages with antigenic novelty or
immune-escape properties could acquire a reproductive advantage even
when their success did not arise solely from increased intrinsic
transmissibility \citep{Carabelli2023}. In the language of the present
framework, prior circulation did several things at once:

\begin{enumerate}
\def\labelenumi{\arabic{enumi}.}
\tightlist
\item
  it created a very large population of already human-adapted viruses;
\item
  it generated standing genetic diversity and continued exploratory
  opportunity through further replication;
\item
  it changed the environment \(E\) by altering host immunity;
\item
  it changed which parts of successor space had high retention or
  establishment probability \(\pi_{\rm est}\).
\end{enumerate}

A lineage approaching its effective persistence boundary could therefore
be replaced by a nearby descendant that restored \(R_e>1\) under the new
immune environment.

This is an eco-evolutionary feedback: realized organization changes the
environment that subsequently changes the accessibility of its
descendants \citep{Ferriere2013}. It is offered only as an illustration
of the general mechanism. The toy model of Sections 2--6 is not a
mechanistic model of SARS-CoV-2 evolution.

\subsection{The inheritance corollary: search does not restart from
zero}\label{the-inheritance-corollary-search-does-not-restart-from-zero}

Exploratory opportunity is only half of Mechanism II. The other half is
the location from which exploration begins.

Suppose the first appearance of a viable organization requires a
difficult sequence of changes:

\[
0
\longrightarrow
\theta.
\]

Once \(\theta\) is reproduced with sufficient fidelity, later variants
are generated conditionally on that realized organization:

\[
\theta
\longrightarrow
\theta+\delta\theta.
\]

The second search is therefore not a repeat of the first. Capabilities
already embodied in \(\theta\) are typically inherited unless subsequent
variation disrupts them. In a pathogen, descendants of a successfully
transmitting lineage usually begin with the molecular and physiological
machinery that already permits host-cell entry, replication, shedding,
and transmission. In the \(A\)--\(B\)--\(C\) example, descendants of the
established architecture begin with the return path already present
unless a variant removes it.

The important claim is therefore not that every descendant preserves
every previous function. It is that the distribution
\(q(\theta'\mid\theta)\) is conditioned on the realized organization.
Its support is centred on transformations of \(\theta\), not on
independent redraws from an organization-free state.

This is the inheritance asymmetry:

\[
\boxed{
\text{successful previous organization becomes the starting point for the next search.}
}
\]

The original establishment problem may have required crossing a large
region of organizational state space. Subsequent adaptive change can be
local because previous solutions are carried into the initial condition
of the next search.

In terms of organizational accessibility, inheritance changes the
conditional distribution of successors even if the underlying mutation
or variation rules are unchanged:

\[
q(\theta'\mid\theta_1)
\neq
q(\theta'\mid\theta_0)
\]

simply because \(\theta_1\neq\theta_0\).

This should be distinguished from evolvability. Inheritance changes the
\textbf{conditioning state} of the variation process. Evolvability,
developed in Section 8, changes the \textbf{variation operator itself}.

\subsection{Organizational accessibility unifies the
mechanisms}\label{organizational-accessibility-unifies-the-mechanisms}

Mechanisms I and II are not two mathematical forms of the same
bifurcation and should not share the label \emph{hysteresis}. They are
two ways in which realized organization changes the general
accessibility kernel defined in Eq. (1). Equation (3) quantifies only
the novel-successor search component relevant especially to Mechanisms
II and III.

\textbf{Mechanism I --- hysteretic retention.} The realized organization
changes the persistence component of \(\mathcal A\). The same external
environment can retain an organization that could not establish there
from rarity. What changes is \textbf{what can remain}.

\textbf{Mechanism II --- inherited adaptive extension.} The realized
organization changes the conditioning state \(\theta\) from which
variants are generated, and its reproductive history changes exploratory
opportunity \(\Omega_\tau\). What changes is \textbf{where search begins
and how much local search occurs}.

A third mechanism will alter the variation process \(Q\) itself, through
\(q\), the event rate \(\nu\), or both.

\begin{longtable}[]{@{}
  >{\raggedright\arraybackslash}p{(\columnwidth - 6\tabcolsep) * \real{0.2500}}
  >{\raggedright\arraybackslash}p{(\columnwidth - 6\tabcolsep) * \real{0.2500}}
  >{\raggedright\arraybackslash}p{(\columnwidth - 6\tabcolsep) * \real{0.2500}}
  >{\raggedright\arraybackslash}p{(\columnwidth - 6\tabcolsep) * \real{0.2500}}@{}}
\toprule\noalign{}
\begin{minipage}[b]{\linewidth}\raggedright
Mechanism
\end{minipage} & \begin{minipage}[b]{\linewidth}\raggedright
Component of accessibility affected
\end{minipage} & \begin{minipage}[b]{\linewidth}\raggedright
What changes
\end{minipage} & \begin{minipage}[b]{\linewidth}\raggedright
Representative formal object
\end{minipage} \\
\midrule\noalign{}
\endhead
\bottomrule\noalign{}
\endlastfoot
I --- hysteretic retention & persistence mass of \(\mathcal A_\tau\) &
what can remain & \(J_{\rm fold}<J_{\rm inv}\) \\
II --- inherited adaptive extension & \(\theta\) and \(\Omega_\tau\) &
where search begins and how many trials occur & \(q(\theta'\mid\theta)\)
and \(\Omega_\tau\simeq \nu N_0/(1-R)\) \\
III --- evolvability & variation process \(Q_m=(q_m,\nu_m)\) & how
candidate successors are generated & \(q_m\) and/or \(\nu_m\) \\
\end{longtable}

The mechanisms can compound. A feedback established through Mechanism I
may preserve organization under conditions that would otherwise
eliminate it. Continued reproduction can then provide the inherited base
point and exploratory opportunity required for Mechanism II. A
successful successor becomes a new base point for subsequent
accessibility:

\[
\theta_t
\rightarrow
\mathcal A_t
\rightarrow
\theta_{t+1}
\rightarrow
\mathcal A_{t+1}
\rightarrow\cdots
\]

No teleology is implied. The system does not estimate \(\varepsilon\),
seek a rescuing value of \(s\), or attempt to maximize \(\mathcal A\).
Variation is generated by the ordinary processes represented by \(Q\)
and \(\Omega\); the environment and current organization determine which
variants establish and how rapidly lineages decline or expand. Near
threshold, the consequences of small differences in reproductive
performance become large because the dynamical filter changes, not
because the system acquires a purpose.

Nor must accessibility widen. A realized organization can narrow future
possibilities, trap itself on a local configuration, or fail to generate
a rescuing successor before exploratory opportunity is exhausted. The
ratchet relevant to cumulative emergence is therefore not an automatic
tendency toward improvement. It is the historical carrying-forward of
organization into the conditions from which later organization is
generated.

\section{Mechanism III --- evolvability: changing how successor search
is
generated}\label{mechanism-iii-evolvability-changing-how-successor-search-is-generated}

Mechanisms I and II change effective organizational accessibility
without requiring the variant-generating process itself to change.
Mechanism I acts through state-dependent persistence and basin
structure. Mechanism II acts through inherited starting architecture and
history-dependent exploratory opportunity. A further recursion becomes
possible when the \textbf{generator of candidate successors itself
varies}.

Let \(m\) denote a heritable or otherwise retainable modifier of the
variation process. Write

\[
Q_m=(q_m,\nu_m),
\]

where \(q_m(\theta'\mid\theta)\) is the conditional distribution of
candidate successor architectures and \(\nu_m\) is their generation rate
per active realization. A modifier may therefore change the
\textbf{shape} of variation, the \textbf{rate} of variation, or both.
Examples include mutation rate, recombination, developmental
variability, exploratory connectivity, modularity, robustness,
genotype--phenotype mapping, the probability that new feedback paths
form, and mechanisms that preserve useful substructure while allowing
other parts to vary.

The modifier changes exploratory opportunity through

\[
\Omega_{m,\tau}
=
\mathbb E\!\left[
\int_t^{t+\tau}N_m(s)\nu_m(s)\,ds
\right],
\]

and hence the novel-successor intensity

\[
\Lambda_{m,\tau}(U)
=
\Omega_{m,\tau}
\int_U
q_m(\theta'\mid\theta)
\pi_{\rm est}(\theta'\mid\theta,z,E)
\,d\theta'.
\]

This decomposition separates two causal routes that can otherwise look
similar. Mechanism II can increase exploratory opportunity because
establishment and subsequent history increase \(N(t)\) or prolong the
time during which the lineage remains extant. Mechanism III can change
exploratory opportunity because the modifier changes \(\nu_m\), and can
independently change the distribution of candidates through \(q_m\).

To connect successor accessibility to persistence consequences, let
\(\Phi(\theta)\) be a persistence-relevant quantity: for example the
range of environments over which an organization remains stable,
expected persistence under an environment distribution, or a
multidimensional measure combining replacement capacity with resource
cost. An accessibility-consistent innovation yield over horizon \(\tau\)
is

\[
\boxed{
Y_\tau(m\mid\theta,z,E)
=
\Omega_{m,\tau}
\int
[\Phi(\theta')-\Phi(\theta)]_+
q_m(\theta'\mid\theta)
\pi_{\rm est}(\theta'\mid\theta,z,E)
\,d\theta'
-
c_\tau(m).
}
\tag{19}
\]

Here \(c_\tau(m)\) is the direct cost of maintaining or expressing
modifier \(m\) over the same horizon. Equation (19) is not a biological
law. It is a bookkeeping expression that places the amount of search,
the distribution of candidates, their probability of establishment,
their persistence consequences, and the direct cost of the variation
machinery in one common formalism.

If a modifier that changes \(q_m\), \(\nu_m\), or both is retained
together with the successful organizations it helps generate, another
feedback becomes possible:

\[
\text{variation-generating architecture}
\rightarrow
\text{changed successor accessibility}
\rightarrow
\text{differential persistence or reproduction}
\rightarrow
\text{retention of the variation-generating architecture}.
\]

The system can then change not only its current organization but the
process that generates distributions of future organization.

This is closely related to \textbf{evolvability}: the capacity to
generate heritable, selectable phenotypic variation
\citep{WagnerAltenberg1996, Gerhart2007}. Theory and experiments show
that mutation rate, recombination, robustness, modularity, and
genotype--phenotype architecture can affect evolvability. Explicit
models have demonstrated conditions under which evolvability can itself
be indirectly selected \citep{EarlDeem2004}.

The qualification is essential. Selection for future adaptive capacity
is generally indirect and may be weaker than selection on immediate
performance \citep{DiazArenas2013, Wagner2022}. A trait can affect
evolvability without having been selected \emph{because} of that effect
\citep{Pelabon2025}. Increased variation can be destructive; robustness
can inhibit useful change as well as protect function; modularity can
constrain as well as facilitate; feedback can preserve maladaptive
organization.

Mechanism III therefore does not mean that systems automatically become
better at adapting. The narrower claim is:

\begin{quote}
\textbf{Self-maintaining systems can, when variant-generating mechanisms
themselves vary and remain coupled to persistence or reproduction,
evolve changes in the rate or distribution by which successor
organizations become accessible.}
\end{quote}

The resulting hierarchy is:

\begin{itemize}
\tightlist
\item
  \textbf{Level 1:} processes produce processes.
\item
  \textbf{Level 2:} a network collectively maintains its organization,
  governed here by \(R_{\rm org}\).
\item
  \textbf{Level 3 --- Mechanism I:} realized feedback changes
  persistence accessibility through basin structure and hysteresis.
\item
  \textbf{Level 4 --- Mechanism II:} inheritance changes the starting
  architecture \(\theta\), while reproductive history changes the
  \(N(t)\) contribution to exploratory opportunity \(\Omega_\tau\).
\item
  \textbf{Level 5 --- Mechanism III:} the organization changes the
  variant-generating process \(Q_m=(q_m,\nu_m)\).
\item
  \textbf{Level 6:} those variation-generating mechanisms themselves
  vary and are retained differentially.
\end{itemize}

At the upper levels, the system can change the process by which
distributions of its own future states are generated. This is a
plausible route to cumulative evolvability, not a demonstrated
consequence of the lower-level model. Modifier dynamics must be modelled
and tested explicitly before any claim of evolved evolvability is
warranted.

\section{Empirical signatures of the three
mechanisms}\label{empirical-signatures-of-the-three-mechanisms}

The framework is useful only if its mechanisms can be distinguished
empirically. The first mechanism yields spectral signatures near
bifurcation thresholds; the second yields a statistical dependence of
rescue on exploratory opportunity and distance from threshold; the third
requires comparison of variant-generating operators.

\subsection{Spectral signatures of the baseline and hysteretic
thresholds}\label{spectral-signatures-of-the-baseline-and-hysteretic-thresholds}

Near a stable equilibrium \(z^*\), add weak stochastic forcing and
linearize:

\[
\dot{\delta z}=A(J)\delta z+\sigma\xi(t).
\]

Projecting onto the slow eigenmode with eigenvalue
\(\lambda_{\rm slow}(J)<0\) yields an Ornstein--Uhlenbeck approximation.
Under standard assumptions,

\[
{\rm Var}(\delta z)\propto\frac{1}{|\lambda_{\rm slow}|},
\qquad
\rho(\Delta t)\approx1-|\lambda_{\rm slow}|\Delta t.
\]

Thus, as the dominant restoring eigenvalue approaches zero, recovery
slows, variance increases, and lagged autocorrelation rises: the
standard critical-slowing-down picture \citep{Scheffer2009, Dakos2012}.

For the baseline transcritical self-maintenance threshold of Section 3,

\[
|\lambda_{\rm slow}|
\propto
|R_{\rm org}-1|,
\qquad
{\rm Var}
\propto
|R_{\rm org}-1|^{-1}.
\]

For the Mechanism-I \(C\)-containing branch approaching its
hysteretic-loss fold from above,

\[
|\lambda_{\rm slow}|
\propto
\sqrt{J-J_{\rm fold}}.
\]

Numerical evaluation of the stable branch gives a fitted exponent of
approximately \(0.515\), consistent with the generic square-root scaling
of a saddle-node. Hence

\[
{\rm Var}
\propto
(J-J_{\rm fold})^{-1/2}.
\]

Two qualifications are important. First, the stable branch must be used
for the return-to-equilibrium interpretation; the opposite side of the
fold is unstable. Second, \(1/\sqrt{\Delta J}\) diverges less strongly
than \(1/\Delta J\) as \(\Delta J\to0\). The fold therefore does not
imply a stronger variance divergence at the same numerical distance from
threshold. What differs is the local scaling law and the geometry of
approach to the transition.

\subsection{Mechanism II: rescue depends on exploratory opportunity and
distance from
threshold}\label{mechanism-ii-rescue-depends-on-exploratory-opportunity-and-distance-from-threshold}

Mechanism II predicts a different signature. For matched systems
beginning decline with the same \(N_0\) and variant-generation rate
\(\nu\), Eq. (15) gives

\[
\Omega
\approx
\frac{\nu N_0}{1-R_{\rm res}}
\]

for a simple subcritical branching approximation. At the same time, Eq.
(17) shows that the minimum relative reproductive advantage required for
rescue is

\[
s_{\min}
=
\frac{1-R_{\rm res}}{R_{\rm res}}.
\]

Pushing the resident further below threshold therefore acts in two
directions simultaneously: it reduces the number of exploratory trials
available before extinction and increases the adaptive displacement
required for any one trial to succeed.

The resulting prediction is not that a ``standing population'' scales as
\(1/(1-R)\). It is that, conditional on comparable initial abundance and
variation machinery, \textbf{rescue probability should decline as
exploratory opportunity falls and the required adaptive step grows}.
Systems with identical current \(R_{\rm res}\) but different histories
may also differ if one enters decline with larger \(N_0\), more standing
variation, or a different \(q\).

A direct test would therefore factorially manipulate:

\begin{enumerate}
\def\labelenumi{\arabic{enumi}.}
\tightlist
\item
  distance below the persistence threshold;
\item
  abundance \(N_0\) at the onset of deterioration;
\item
  standing variation present at that moment;
\item
  the rate or distribution of newly generated variants.
\end{enumerate}

If recovery depends only on the current environment and not on these
history-dependent sources of accessibility, Mechanism II is not
supported.

\subsection{Mechanism III: accessibility changes when the variation
operator
changes}\label{mechanism-iii-accessibility-changes-when-the-variation-operator-changes}

Mechanism III requires a different experiment again. Compare otherwise
matched organizations with modifiers \(m_1\) and \(m_2\) that differ in
\(q_m\), \(\nu_m\), or both, while controlling as far as possible for
direct present-day effects.

The relevant observations are not simply whether one lineage adapts
faster, but whether the modifiers measurably change:

\begin{itemize}
\tightlist
\item
  the distribution of candidate successors generated;
\item
  the probability that generated candidates occupy persistence-improving
  regions of state space;
\item
  the number of candidate-generation events per unit reproductive
  exposure;
\item
  and, crucially, whether the modifier is differentially retained
  because of those consequences.
\end{itemize}

The final requirement is what distinguishes \textbf{selection on
evolvability} from a trait that merely happens to alter evolvability as
a side effect.

All of these predictions should be treated cautiously. Empirical
early-warning signals depend on noise structure, sampling,
nonstationarity, rates of parameter change, finite population size, and
observability of the relevant state variables. The equations predict
signatures of the present model and its simple stochastic extensions;
they do not establish that any particular biological, social,
technological, or artificial system is governed by the same local normal
forms.

\section{What the framework explains --- and what it does
not}\label{what-the-framework-explains-and-what-it-does-not}

The construction supports six claims.

\begin{enumerate}
\def\labelenumi{\arabic{enumi}.}
\item
  \textbf{Collective self-maintenance is possible without a central
  controller.} Local production processes can jointly replace the
  processes that execute the network.
\item
  \textbf{Self-maintenance has a threshold.} The base model has a
  reproduction number \(R_{\rm org}\) that separates decay from
  establishment at low abundance.
\item
  \textbf{Realized feedback can change its own retention conditions.} A
  new loop can produce bistability and hysteresis, so that an
  organization persists under conditions in which it could not establish
  from rarity. This is Mechanism I.
\item
  \textbf{Inheritance and reproductive history can change successor
  accessibility without hysteresis.} Once an organization has been
  realized, descendants are generated around that already viable
  structure, and the number of exploratory trials available before loss
  depends on abundance, reproduction, and variant-generation rate. Near
  a persistence boundary, the required adaptive displacement may be
  small while exploratory opportunity remains substantial. This is
  Mechanism II.
\item
  \textbf{History can change effective organizational accessibility, not
  merely the currently occupied state.} A realized successor can change
  basin occupancy, inherited architecture \(\theta\), exploratory
  opportunity \(\Omega\), the future environment, or the
  variant-generating process \(Q\). Consequently, the accessibility
  encountered later in the lineage may differ because of the
  organization previously realized.
\item
  \textbf{There is a coherent route to evolvability feedback.} If
  variant-generating mechanisms themselves vary and remain coupled to
  the persistence or reproduction of the organizations they help
  produce, selection can act on changes in \(q_m\), \(\nu_m\), or
  related innovation machinery. This is Mechanism III.
\end{enumerate}

The framework does \textbf{not} show that:

\begin{itemize}
\tightlist
\item
  complexity must increase monotonically;
\item
  accessibility must widen over time;
\item
  every added process or connection is useful;
\item
  every persistent organization is beneficial by an external criterion;
\item
  more variation, connectivity, or feedback is always advantageous;
\item
  proximity to \(R=1\) guarantees evolutionary rescue;
\item
  cumulative lineage size in Eq. (14) is the same thing as instantaneous
  standing population size;
\item
  inheritance preserves every previously acquired function;
\item
  Mechanism II is a form of hysteresis;
\item
  evolvability is always selected because of its effects on future
  variation;
\item
  the present toy reactions are realized in any particular organism,
  ecosystem, institution, robot, or AI system;
\item
  the number of process types is a sufficient measure of biological or
  informational complexity.
\end{itemize}

The same formalism allows simplification, loss, trapping, and
extinction. A fast self-reproducing competitor can displace a
multi-process network. An inherited architecture can constrain future
search. A lineage pushed far below threshold may still fail to recover
despite a large prior population if its remaining exploratory
opportunity and accessible distribution of variant effects are
insufficient for rescue. A lineage can exhaust its exploratory
opportunity without finding a rescuing successor. A variation modifier
can raise the production of harmful variants.

The framework therefore supplies mechanisms for \textbf{historical
transformation of organizational accessibility}. Cumulative emergence is
one possible regime of that transformation, not its inevitable
direction.

\section{A research programme for organizational
accessibility}\label{a-research-programme-for-organizational-accessibility}

The framework becomes scientifically useful when its components are
measured separately.

For a candidate real system, one should identify:

\begin{enumerate}
\def\labelenumi{\alph{enumi})}
\item
  the organizational architecture \(\theta\): the process types,
  interactions, or other causal relations whose continuation defines the
  organization;
\item
  the resources consumed in their replacement;
\item
  the directed causal effects by which one process increases the future
  abundance or execution of another;
\item
  the reproduction, replacement, or invasion threshold of the existing
  organization;
\item
  a candidate innovation and its opportunity cost;
\item
  whether the innovation creates a return path into the organization
  that produces or maintains it;
\item
  whether establishment and maintenance thresholds differ, indicating
  Mechanism I;
\item
  whether current persistence under identical external conditions
  depends on historical seeding or initial conditions;
\item
  the dynamical state \(z(t)\) and, where relevant, the active
  population or number of realizations \(N(t)\) through time;
\item
  the rate \(\nu(t)\) at which relevant candidate successors are
  generated, allowing estimation of exploratory opportunity
\end{enumerate}

\[
\Omega_\tau=\int_t^{t+\tau}N(s)\nu(s)\,ds;
\]

\begin{enumerate}
\def\labelenumi{\alph{enumi})}
\setcounter{enumi}{10}
\item
  the distribution \(q(\theta'\mid\theta)\) of candidate successors
  around the current organization;
\item
  the establishment filter \(\pi_{\rm est}(\theta'\mid\theta,E)\) for
  generated candidates;
\item
  whether recovery from deterioration depends on the inherited starting
  state, abundance at deterioration onset, standing variation, and
  \(\Omega\), as predicted by Mechanism II;
\item
  whether variation-generating mechanisms differ among persistent
  lineages or organizations;
\item
  whether such modifiers alter \(q\), \(\nu\), or both;
\item
  whether those modifiers are retained because of the successor
  organizations they make accessible rather than because of unrelated
  direct benefits.
\end{enumerate}

The decisive experiments differ by mechanism.

For \textbf{Mechanism I}, intervene on return paths. Introduce the same
candidate process in otherwise matched systems, preserve its resource
cost, and selectively break the proposed feedback from consequence to
renewed production. If the difference between establishment and
maintenance remains unchanged, the proposed loop was not causal.

For \textbf{Mechanism II}, manipulate accessibility without changing the
retention architecture. Matched systems can be driven below threshold
with different \(N_0\), different amounts of standing variation, or
different prior reproductive histories while holding the current
environment fixed. The prediction is that recovery depends on the
resulting \(\Omega\) and on the distribution of local variants, not only
on current \(R_{\rm res}\).

A particularly informative experiment would compare the accessibility of
a target viable organization from two starting points:

\[
0\longrightarrow\theta'
\]

versus

\[
\theta\longrightarrow\theta',
\]

where \(\theta\) already contains part of the causal organization
required for viability. The inheritance claim predicts that the second
transition can be substantially more accessible when prior viable
structure is preserved.

For \textbf{Mechanism III}, compare architectures differing in their
variant-generating operator while controlling direct present-day
performance. Measure not merely the number of variants produced but the
resulting distribution over organizational state space and the fraction
surviving the establishment filter.

The central empirical object is therefore not ``complexity'' in the
abstract. It is the causal decomposition of accessibility:

\[
\boxed{
\text{starting organization}
\times
\text{exploratory opportunity}
\times
\text{variation distribution}
\times
\text{retention dynamics}.
}
\]

The factorization is conceptual rather than universally multiplicative,
but it identifies quantities that can be perturbed and measured
separately.

\section{Discussion: from persistence to the historical transformation
of
accessibility}\label{discussion-from-persistence-to-the-historical-transformation-of-accessibility}

Darwinian natural selection explains differential propagation of
heritable variation. The present framework addresses a complementary
question: \textbf{what determines the organizational variation and
persistent states that are effectively accessible for selection and
dynamics to act upon?} The answer proposed here is historical. Realized
organization can become part of the causal conditions determining what
can remain, where successor search begins, how much search occurs, and
how candidate successors are generated.

The distinction from ordinary state dependence matters. A fixed
dynamical law may of course give different futures from different
current states. The present theory does not need to claim that every
organizational transition rewrites that law. Its claim is instead that
organizational history can systematically alter \textbf{effective
accessibility} through several mechanistically distinguishable channels.
Hysteretic basin occupancy can retain a state at fixed external
conditions. Inheritance can make previously viable architecture the base
point of later variation. Reproductive history can alter exploratory
opportunity. Eco-evolutionary feedback can alter future environment.
Evolvability can change the generator \(Q\) itself.

Mechanism I provides the most explicit dynamical demonstration. In the
seed-dependent model, the conditions for establishment and maintenance
separate: \(J_{\rm fold}<J_{\rm inv}\). The same external environment
can therefore support different persistent outcomes depending on which
basin history has placed the system in. This is genuine hysteresis and
should remain the only mechanism in the paper described by that term.

Mechanism II is different. Once a viable architecture is inherited,
successor variation is conditioned on that architecture rather than
drawn independently from an organization-free state. Continued activity
also supplies a finite exploratory opportunity \(\Omega_\tau\). Near a
persistence boundary, the same lineage may simultaneously decline slowly
enough to generate many remaining trials and require only a small
favourable displacement to restore positive growth. That combination is
the setting of evolutionary rescue, but the deeper inheritance result
does not depend on rescue: \textbf{successful previous organization
becomes the starting point for the next search}.

Mechanism III takes the recursion one level further. If organizations
differ in \(Q_m=(q_m,\nu_m)\), then they differ not merely in present
performance but in the distribution and rate of future variants they
generate. Selection can then, under restricted conditions, act
indirectly on properties of the search process itself. Equation (19)
makes this connection explicit by combining exploratory opportunity,
candidate distribution, establishment, persistence consequences, and the
cost of the modifier.

These mechanisms do not imply progress. A feedback can preserve a
maladaptive state. Inheritance can trap a lineage near a local
configuration. Exploratory opportunity generates harmful as well as
useful candidates. A changed variation operator can reduce rather than
increase access to viable successors. Evolution may simplify
organization or terminate lineages. The general statement is therefore
not that existence makes subsequent existence easier. It is:

\[
\boxed{\text{Existence changes what existence can follow.}}
\]

The cumulative case requires additional conditions. If organization
retains causal structure needed for viability, if that structure is
inherited with sufficient fidelity, and if subsequent variation can
modify it without reconstructing it from zero, then previous viability
can become infrastructure for later viability:

\[
\boxed{\text{previous viability can become infrastructure for subsequent viability}.}
\]

This provides a precise sense in which organizational history can
accumulate without invoking a force toward complexity. What accumulates
need not be complexity as such; it is \textbf{causal structure that
conditions later accessibility}.

The relationship to Darwinian selection can then be stated simply.
Selection does not operate on all logically imaginable forms. It acts on
variants that are generated, encountered, and able to establish under
particular ecological and dynamical conditions. The accessibility
framework asks why that presented distribution has the structure it
does. Previous organization may have altered the inherited base point,
the amount of search, the establishment landscape, the environment, or
the variation generator itself. In that sense:

\begin{quote}
\textbf{Selection acts within an accessibility structure that previous
organization may itself have helped create.}
\end{quote}

The recursion is therefore better written with the full conditioning
state than as a claim that a mathematical kernel is literally rewritten
at every step:

\[
(\theta_t,z_t,E_t,Q_t)
\longrightarrow
\mathcal A_t
\longrightarrow
(\theta_{t+1},z_{t+1},E_{t+1},Q_{t+1})
\longrightarrow
\mathcal A_{t+1}
\longrightarrow\cdots
\]

The important historical possibility is that

\[
\mathcal A_{t+1}\neq\mathcal A_t
\]

because one or more of the causal conditions determining accessibility
were changed by what happened in the intervening transition. Only
Mechanism III requires the generator \(Q\) itself to change; Mechanisms
I and II can transform effective accessibility through state,
inheritance, and history while the underlying variation rules remain
unchanged.

This suggests a restricted use of the term \emph{ratchet}:

\begin{quote}
\textbf{A cumulative organizational ratchet occurs when realized
organization changes subsequent effective accessibility in a way that
carries forward enough previously achieved viability that later viable
organization need not be rediscovered de novo.}
\end{quote}

Mechanism I can carry viability forward by making an established state
harder to lose than it was to establish. Mechanism II can carry it
forward because descendants start from inherited viable architecture and
retain a finite, history-dependent opportunity to explore nearby
alternatives. Mechanism III can additionally carry forward properties of
the process by which alternatives are generated.

The deepest empirical question is therefore not whether systems tend
toward complexity, but:

\begin{quote}
\textbf{When, and by what mechanisms, does realized organization become
part of the causal conditions that determine future organizational
accessibility?}
\end{quote}

The present models provide one constructive answer. They do not
establish a universal law.

\section{Conclusion}\label{conclusion}

This paper develops a dynamical account of \textbf{organizational
accessibility}: the finite-horizon accessibility of successor
organizational states conditional on organizational architecture
\(\theta\), dynamical state \(z\), external environment \(E\), and
variant-generating process \(Q=(q,\nu)\).

The argument has four layers.

First, an open network of local production processes can collectively
cause its own continuation. Its establishment is governed by the
reproduction threshold \(R_{\rm org}\). This baseline determines whether
organization can grow from low abundance at all.

Second, \textbf{Mechanism I --- hysteretic retention} shows how realized
state can change what remains accessible at fixed external conditions.
In the worked example, a seed-dependent innovation invades only above
\(J_{\rm inv}=1.0667\) but, once established, persists down to
\(J_{\rm fold}=0.5978\). The resulting backward bifurcation creates
bistability and history-dependent basin occupancy.

Third, \textbf{Mechanism II --- inherited adaptive extension} shows how
realization changes where successor search begins and how much search
can occur. Reproduction with inheritance makes already viable
architecture the base point of later variation. Reproductive history
generates exploratory opportunity

\[
\Omega_\tau
=
\mathbb E\!\left[
\int_t^{t+\tau}N(s)\nu(s)\,ds
\right].
\]

For a simple subcritical branching process, cumulative lineage size per
founder scales as \(1/(1-R)\) as \(R\to1^-\); this is a statement about
cumulative reproductive opportunity, not instantaneous standing
population. Near the same boundary, the adaptive displacement required
to restore positive growth may be small. Together these effects create
the conditions for evolutionary rescue without requiring hysteresis.

Fourth, \textbf{Mechanism III --- evolvability} extends the recursion to
the variant-generating process itself. If organizations differ in
\(Q_m=(q_m,\nu_m)\) and those differences remain coupled to descendant
persistence or reproduction, then selection can alter how candidate
successors are generated. The search process becomes an evolvable
component of organizational accessibility.

The common formal object is

\[
\boxed{
\mathcal A_\tau(S\mid\theta,z,E,Q)
=
\Pr\!\left(
(\Theta_{t+\tau},Z_{t+\tau})\in S
\mid
\Theta_t=\theta,
Z_t=z,
E,Q
\right).
}
\]

For novel successors, the corresponding successful-successor intensity
is

\[
\boxed{
\Lambda_\tau(U)
=
\Omega_\tau
\int_U
q(\theta'\mid\theta)
\pi_{\rm est}(\theta'\mid\theta,z,E)
\,d\theta'.
}
\]

This decomposition distinguishes \textbf{how many opportunities occur},
\textbf{what variants are generated}, and \textbf{which generated
variants can establish}. Inheritance determines the architecture from
which that distribution is generated; hysteresis determines whether
already-realized organization remains accessible; evolvability changes
the generator itself.

The framework does not imply that accessibility increases, that
complexity accumulates monotonically, or that evolution possesses
foresight. History may narrow accessibility, preserve liabilities,
constrain search, or end in simplification and extinction.

The general claim is instead:

\[
\boxed{\text{Existence changes what existence can follow.}}
\]

Under conditions that preserve enough causal structure for viability,
the cumulative special case is:

\[
\boxed{\text{What persists becomes part of the starting conditions for what can persist next.}}
\]

Equivalently, previous viability can become infrastructure for
subsequent viability. Cumulative emergence is therefore not
fundamentally the accumulation of complexity. It is a historical regime
in which realized organization conditions future accessibility while
carrying forward enough of what already worked that later viability need
not repeatedly be discovered from zero.

Selection then operates within an accessibility structure that previous
organization may itself have helped create. That is the proposed
dynamical basis of cumulative organization.


\printbibliography

\end{document}
